% ============================================================================
%  Task_W01_Answers.tex -- Fillable answer sheet for Task_W01: Recursion Basics
%  Uses coa-accessible-student.sty for a consistent look with the
%  assignment itself, without the AI-development disclaimer.
%
%  Students: see the "Before You Start" section in the compiled PDF (right
%  after the Name line below) for how to edit and compile this file.
% ============================================================================
\documentclass[11pt]{article}
\usepackage{coa-accessible-student}

% Same plain pseudocode listing style used elsewhere in the course.
\lstdefinestyle{pseudo}{
  basicstyle=\ttfamily\small\color{coaBodyText},
  frame=single,
  rulecolor=\color{coaNavy},
  breaklines=true,
  showstringspaces=false,
  columns=fullflexible,
  tabsize=2,
  numbers=none
}
\lstset{style=pseudo}

\title{Task W01: Recursion Basics --- Answer Sheet}
\hypersetup{pdftitle={Task W01: Recursion Basics -- Answer Sheet}, pdflang={en-US}}

\begin{document}

{\centering
{\Large\bfseries\color{coaNavy} Task W01: Recursion Basics --- Answer Sheet\par}
\vspace{0.6em}
}
\noindent Name: \underline{\hspace{3.2in}}
\vspace{1em}

\begin{center}
\rule{0.95\linewidth}{0.6pt}
\end{center}
\subsection*{Before You Start}
\noindent
This is a LaTeX source file, not a finished document; what you are reading
right now is what it looks like once compiled. To edit it and produce your
own PDF, you will need a LaTeX environment such as Overleaf, TeXworks, or a
local TeX installation; if you need help setting it up in Linux, please let
me know.

\begin{enumerate}[leftmargin=1.4em]
  \item Keep \texttt{Task\_W01\_Answers.tex} and
        \texttt{coa-accessible-student.sty} in the same folder. This file
        will not compile without the style file present.
  \item Open \texttt{Task\_W01\_Answers.tex} in a text or code editor, not a
        word processor, and fill in each section below with your own work,
        following the instructions given in each section.
  \item Compile it by running \texttt{pdflatex Task\_W01\_Answers.tex}. Run
        this command twice in a row; the page numbers at the bottom of each
        page will not resolve correctly on a single pass, and will show
        "Page 1 of ??" instead of the real total.
  \item Every time you make a change to the \texttt{.tex} file, save it and
        run \texttt{pdflatex} twice again to see your latest edits in the
        PDF you submit.
  \item Submit both files: your edited \texttt{Task\_W01\_Answers.tex}
        source and the \texttt{Task\_W01\_Answers.pdf} it produces. Your
        diagrams are submitted separately from this file; see the reminder
        in each section below and the Diagram Submission Instructions in
        the assignment itself.
\end{enumerate}
\begin{center}
\rule{0.95\linewidth}{0.6pt}
\end{center}
\vspace{0.6em}

\section{Recurrence Relation Problems}
Diagrams for both problems below are required, but are \textbf{not}
included in this file. Submit them separately, following the Diagram
Submission Instructions in the assignment. Just record your final answer
here.

\subsection{Problem: $B(n)$}
Value of $B(6)$: \underline{\hspace{1.5in}}

\subsection{Problem: $D(n)$}
Value of $D(6)$: \underline{\hspace{1.5in}}

\section{Applied Problems}
Problem you are completing in full (write the number 1 through 4):
\underline{\hspace{1in}}

\noindent Final returned value for that problem: \underline{\hspace{1.5in}}

The diagram for whichever problem you complete in full is required, but is
\textbf{not} included in this file. Submit it separately, following the
Diagram Submission Instructions in the assignment. The other three problems
do not need a diagram, just your attempted solution below.

You may write pseudocode or working code in Java or C/C++ for each problem
below; either is fine, and you do not need to make the same choice for
every problem. If you would like to see it written the way CLRS and Walls
and Mirrors write it, see the link in the assignment itself. If you use the
AI-for-formatting exception described there, submit that conversation as
its own file, the same way you submit your diagrams separately.

\subsection{Problem 1}
\begin{lstlisting}
% Your pseudocode or code here.
\end{lstlisting}

\subsection{Problem 2}
\begin{lstlisting}
% Your pseudocode or code here.
\end{lstlisting}

\subsection{Problem 3}
\begin{lstlisting}
% Your pseudocode or code here.
\end{lstlisting}

\subsection{Problem 4}
\begin{lstlisting}
% Your pseudocode or code here.
\end{lstlisting}

\section{Bugs or Issues You Found}
If you found something in the assignment that needs to be debugged, for
example an unclear problem statement or an example that does not behave
the way it is described, describe it here. Leave this blank if nothing
came up.

\vspace{3em}

\end{document}
